%% =====================================================================
%%  Representation rather than architecture determines cross-cohort
%%  transfer in automated scalp EEG seizure detection
%%
%%  Compile:  pdflatex main ; bibtex main ; pdflatex main ; pdflatex main
%%  On Overleaf set the compiler to pdfLaTeX.
%%
%%  Standard article class, so it builds anywhere. Instructions for
%%  switching to the IOP class at submission are at the end of the file.
%% =====================================================================

\documentclass[11pt,a4paper]{article}

\usepackage[top=2.0cm,bottom=2.0cm,left=2.0cm,right=2.0cm]{geometry}
\usepackage[T1]{fontenc}
\usepackage[utf8]{inputenc}
\usepackage{lmodern}
\usepackage{graphicx}
\usepackage{amsmath,amssymb}

%% Missing figures render as a labelled box rather than breaking the build.
%% Put the PDFs in figures/ and they appear automatically; nothing else to
%% change. Delete this block once every figure is in place.
\graphicspath{{figures/}{./}}
\makeatletter
\let\orig@includegraphics\includegraphics
\renewcommand{\includegraphics}[2][]{%
  \IfFileExists{#2}%
    {\orig@includegraphics[#1]{#2}}%
    {\fbox{\parbox[c][3.8cm][c]{0.86\textwidth}{\centering
      \textsf{\small figure not found}\\[2pt]\texttt{\footnotesize #2}}}}}
\makeatother
\usepackage{booktabs}
\usepackage{siunitx}
\usepackage{microtype}
\usepackage[numbers,sort&compress]{natbib}
\usepackage[hidelinks]{hyperref}
\usepackage{authblk}
\usepackage{titlesec}
\usepackage{caption}


\newcommand{\etal}{\textit{et~al.}}

\captionsetup{font=small,labelfont=bf,skip=5pt}
\titleformat{\section}{\normalfont\large\bfseries}{\thesection.}{0.6em}{}
\titleformat{\subsection}{\normalfont\normalsize\bfseries}{\thesubsection.}{0.6em}{}
\titlespacing*{\section}{0pt}{1.4ex plus .4ex}{0.8ex}
\titlespacing*{\subsection}{0pt}{1.2ex plus .3ex}{0.6ex}
\setlength{\parskip}{0.25em}
\setlength{\parindent}{1.2em}
\setlength{\textfloatsep}{10pt plus 2pt minus 2pt}
\setlength{\intextsep}{10pt plus 2pt minus 2pt}
\linespread{1.04}
\setlength{\bibsep}{1pt}

\renewcommand{\arraystretch}{1.08}
\sisetup{detect-weight=true,detect-family=true}

\title{\vspace{-1.2cm}\bfseries Representation rather than architecture
dominates cross-cohort transfer in automated scalp EEG seizure detection}

\author[1]{Anurag Tiwari}
\affil[1]{\small Department of Computer Science and Engineering, South Asian
University, Maidan Garhi, New Delhi 110068, India\\
\texttt{anuragtiwari@sau.int}}

\date{}

\begin{document}

\maketitle
\vspace{-0.8cm}

\begin{abstract}
\noindent\small
\textbf{Objective.} Automated seizure detection is usually developed and
validated within a single dataset, and the difficulty of transferring a
detector to another recording centre is generally assumed to exceed that of
transferring it to a new patient. We measured both under one protocol, and
asked what determines the difference.
\textbf{Approach.} Fourteen published architectures spanning 2.4$\times$10$^3$
to 1.1$\times$10$^6$ parameters, and two handcrafted spectral feature
families, were evaluated on CHB-MIT (24 recording sets from 23 paediatric
subjects, \SI{256}{\hertz})
and Siena (14 adult subjects, \SI{512}{\hertz}) across three nested settings:
within-patient, cross-patient within a cohort, and cross-cohort in both
directions. All models saw identical windows, splits and preprocessing.
Evaluation prevalence was then swept from balanced to the native
\SI{0.34}{\percent} with the trained model held fixed.
\textbf{Main results.} Learned representations retained substantially more
performance under cohort transfer than handcrafted features, while differing
little among themselves (0.785--0.865 across fourteen architectures). Applied
without retraining to all \num{1000311} windows at the native
\SI{0.34}{\percent} prevalence, they reached a mean cross-cohort AUROC of
0.843 against 0.686, a lift over chance of 42.5$\times$ against
2.8$\times$. Crossing a cohort boundary was consistently no more
costly than crossing a patient boundary, and generally less so, though the
magnitude is imprecisely estimated: in the nested comparison the two drops
were 0.039 and 0.063 AUROC, a ratio of 0.62 (\SI{95}{\percent} CI 0.37 to
1.35), and in the reverse direction 0.32 (0.02 to 0.72). In 13 of 24
leave-one-out folds a new patient from the same centre was harder than the
other cohort entirely.
Cohort identity and subject identity were similarly recoverable from the
features (AUROC 0.92 and 0.95), both exceeding the seizure label itself
(0.81), and the per-feature magnitude of cohort shift was uncorrelated with
that feature's discriminative importance ($r=-0.096$). Area under the
precision--recall curve fell by a factor of 165 across the range of
evaluation prevalences used in the literature while AUROC moved by 0.033.
Across four architectures rerun with five seeds each, the spread between
architectures was 0.66 times the spread from initialisation alone in the
cross-patient setting and 1.09 times it cross-cohort, and which architecture
ranked best cross-cohort changed with the seed. Scored at the event level on continuous recording, the learned
representations detected \SI{22}{\percent} of cross-cohort seizures at one
false alarm per hour and \SI{50}{\percent} at two, against \SI{1}{\percent}
and \SI{17}{\percent} for the handcrafted features.
\textbf{Significance.} Under this protocol, cross-cohort transfer appears to
depend more on the choice of representation than on the choice of
architecture, and more on inter-patient variability than on the additional
cohort shift observed between these two cohorts. A substantial part of the spread
in published performance figures is attributable to evaluation protocol.
\end{abstract}

\vspace{-0.2cm}
\noindent\small\textbf{Keywords:} seizure detection, scalp EEG,
generalization, domain shift, evaluation methodology, benchmarking
\normalsize

\vspace{0.3cm}

%% =====================================================================
\section{Introduction}
%% =====================================================================

Automated seizure detection from scalp EEG has been reported at near-ceiling
performance for more than a decade. On the CHB-MIT Scalp EEG Database
\citep{shoeb2009} accuracies above \SI{98}{\percent} are routine, and several
recent systems report figures above \SI{99}{\percent} together with
false-prediction rates below 0.006 per hour \citep{hybridconv2025}. Clinical
adoption has not followed. The discrepancy is usually attributed to
distribution shift: a detector developed at one centre, on one acquisition
system, in one patient population, is expected to fail at another. That
expectation has motivated a substantial literature on domain adaptation for
EEG, in which the target domain is variously defined as an unseen subject or
an unseen dataset \citep{jemal2024}.

The expectation is plausible but has not been measured
(figure~\ref{fig:boundaries}). Cross-subject and cross-dataset
generalization are treated as related but distinct problems,
and the latter is generally presented as the harder of the two, yet the two
are rarely evaluated side by side on identical windows, under identical
splits, with identical models. Without that comparison it cannot be
established whether a method that improves cross-dataset transfer is
addressing dataset shift specifically, or is addressing inter-subject
variability that is equally present within a single cohort. The distinction
matters practically: if the residual error after aligning two datasets is
inter-subject rather than inter-dataset, then the alignment is operating on
the wrong axis, and the reported improvement would also have been obtained by
pooling more subjects.

A second difficulty is that reported figures are not comparable between
studies. A recent survey of cross-subject EEG decoding identifies the absence
of standardized benchmarking pipelines as a principal obstacle, noting that
differences in artifact rejection thresholds, spectral filtering ranges,
subject exclusion criteria, subject-split strategy and validation protocol
substantially alter task difficulty, so that reported accuracies reflect the
experimental pipeline as much as the intrinsic merit of the method
\citep{crosssubjsurvey2026}. In seizure detection this is acute, because the
positive class is rare: interictal activity dominates continuous recording by
two to three orders of magnitude, and the near-universal practice of
rebalancing before evaluation---by synthetic minority oversampling, by
overlapping-window augmentation, or by majority undersampling---changes the
operating conditions of the reported metric. Saha \etal \citep{saha2024}
examined evaluation practice on CHB-MIT specifically and observed that many
studies do not evaluate on the full recording, and that partial-data
experiments yield high apparent performance that does not survive the
imbalanced continuous setting.

This is an instance of a general property of the metrics involved. Under
random ranking the expected area under the precision--recall curve is the
positive-class rate itself, so that rate is the chance level against which
any value must be read \citep{saito2015}. When it varies by two orders of
magnitude between evaluation protocols, as it does here, a raw value cannot
be interpreted without it. Any metric whose chance level depends on
prevalence will move when prevalence moves, and in seizure detection
prevalence is a design choice rather than a property of the disease. The area under the receiver operating characteristic curve is
prevalence-invariant and so does not exhibit the effect, but it is for that
reason unable to reveal it, and it weights all parts of the ranking equally in
a task where only the top of the ranking is clinically actionable
\citep{saito2015}.

\begin{figure}[t]
\centering
\includegraphics[width=\textwidth]{figures/F7_boundaries.pdf}
\caption{\textbf{The two boundaries and the hypothesis under test.} (a) A
model developed on one cohort can be asked to cross either of two
boundaries: to a different patient within the same cohort, or to a different
cohort entirely. Each ellipse is one patient. (b) Cross-dataset transfer is
generally presented as the harder of the two, which implies a further large
drop at the third point. The size of that drop, relative to the drop already
incurred by meeting a new patient, is what this study measures.}
\label{fig:boundaries}
\end{figure}


We therefore ran a controlled comparison. Two public scalp EEG cohorts were
used that differ in patient age and clinical context as well as in
acquisition hardware, sampling rate, electrode montage and mains frequency,
and that consequently constitute a genuine cohort-level shift rather than a
nominal one. These factors cannot be separated with these data, and the
quantity measured here is the cost of crossing between two cohorts that
differ in all of them at once. Fourteen published
architectures spanning three orders of magnitude in parameter count, together
with two handcrafted spectral feature families, were evaluated under a single
fully specified protocol with three nested settings. Every model saw the same
windows, the same splits and the same preprocessing, so that differences
between settings are attributable to the setting rather than to the pipeline.

We asked three questions. How large is the additional cost of crossing a
dataset boundary, relative to the cost of crossing a patient boundary? Does
that comparison depend on the choice of architecture, or on whether the
representation is learned or hand-designed? And how much of the spread in
reported seizure-detection performance is attributable to evaluation
prevalence rather than to method?

We do not propose a new architecture. The finding of this study is that,
among the architectures examined and at the resolution these data permit,
architecture is not where the limiting factor appears to lie.

%% =====================================================================
\section{Related work}
%% =====================================================================

\subsection{Performance and its interpretation on CHB-MIT}

CHB-MIT \citep{shoeb2009} has been the dominant benchmark for scalp EEG
seizure detection for over a decade, and reported performance has risen
steadily. Recent work using classical feature engineering with explicit
leakage control reports an area under the curve of 0.927 from a compact
fifteen-feature set under strictly patient-independent nested
cross-validation, with feature selection and spatial filtering performed
within training folds only \citep{discoverappl2026}. Deep architectures
routinely report higher figures. The gap between these two families of
results is not straightforwardly attributable to representational capacity,
because the evaluation protocols differ in ways that affect the achievable
value of the metric.

The reliability of the higher figures has been questioned from within the
field. Saha \etal \citep{saha2024} note that studies frequently evaluate on a
subset rather than the full continuous recording. Independently, the same
concern appears in the general EEG decoding literature
\citep{crosssubjsurvey2026}. Our contribution is to quantify the size of the
effect for this task rather than to restate its existence.

\subsection{Cross-subject and cross-dataset transfer}

Cross-patient detectors have long been recognised as substantially harder
than patient-specific ones, since seizure morphology, localization and
duration vary markedly between individuals. Domain adaptation has been the
dominant response. Jemal \etal \citep{jemal2024} evaluated three domain
adaptation methods on CHB-MIT and Siena and reported accuracy improvements of
\SI{10.3}{\percent} and \SI{7.4}{\percent} respectively. Figures of this
kind are difficult to place against one another, because the adapted and
unadapted baselines are rarely evaluated under a common protocol: an
improvement reported over a weak baseline and an absolute figure obtained
under a different windowing or class balance are not on the same scale. The
comparison made here---identical windows, splits and preprocessing across
every model and every setting---is intended to remove that particular source
of ambiguity.

Cross-dataset evaluation between CHB-MIT and Siena has been performed in both
directions and in a pooled configuration \citep{hybridconv2025}, with $F_1$
exceeding \SI{95}{\percent} reported for pure cross-dataset transfer. That
evaluation follows a pipeline in which the minority class is expanded by
synthetic minority oversampling and by sliding windows with \SI{80}{\percent}
overlap, and the majority class is then undersampled to match.
Overlapping-window augmentation at this rate produces near-duplicate windows
within a single seizure episode, and rebalancing prior to evaluation fixes the
prevalence at which the metric is computed. Neither choice is incorrect in
itself, and both are declared; the difficulty is that a figure obtained under
those conditions and a figure obtained on continuous recording are not
comparable, and nothing in the reported metrics indicates by how much they
differ.

\subsection{Standardized benchmarking and null models}

The need for standardized evaluation has been recognised across EEG
decoding more broadly, and the survey cited above identifies the absence of
common pipelines as a principal obstacle to comparing reported figures
\citep{crosssubjsurvey2026}. That literature addresses breadth: common
interfaces across many tasks and datasets. The present study is
complementary in addressing depth. We hold one clinical task fixed and vary
the generalization setting and the evaluation prevalence, so that the
contribution of protocol to reported performance can be isolated from the
contribution of method.

A directly analogous comparison has been reported outside seizure detection.
In consumer-grade EEG emotion recognition, a systematic comparison of EEGNet
against traditional feature engineering with cross-dataset validation found
that handcrafted features outperformed end-to-end learning in that
low-channel, high-noise regime \citep{consumereeg2026}. Our result runs in the
opposite direction, which is informative rather than contradictory: the
advantage of learned representations appears to be regime-dependent, and a
comparison of the two regimes constrains when either should be preferred.

The principle underlying our controls is established. Conrad \etal
\citep{conrad2022} constructed a spatial null that predicts which intracranial
electrodes lie in the seizure onset zone from electrode placement alone, with
no EEG data, and found that it classified electrodes at an area under the
curve of 0.70, and that adding functional connectivity did not significantly
improve on it; they argued that a feature should be required to outperform
such a null before being called a biomarker. The present study applies that
requirement in two directions at once: to the metric, by asking what a
reported figure owes to the prevalence at which it was computed, and to the
comparison between datasets, by asking what the difference between two
reported figures owes to the conditions under which each was measured.

%% =====================================================================
\section{Methods}
%% =====================================================================

\subsection{Datasets}

\textbf{CHB-MIT Scalp EEG Database} \citep{shoeb2009}. Paediatric subjects
with intractable focal epilepsy recorded at Children's Hospital Boston during
pre-surgical evaluation, with antiseizure medication tapered. Recordings are
continuous at \SI{256}{\hertz} in a bipolar longitudinal montage, with seizure
onsets and offsets annotated in per-subject summary files.

The database contains 24 recording sets from 23 individuals: chb21 is the
same patient as chb01, recorded some eighteen months later. We treat the two
as separate units throughout, as the database itself does, and report counts
as 24 for that reason. This is a deliberate choice with a cost: a
cross-patient split that places chb01 and chb21 on opposite sides is not
strictly patient-independent. Their windows are separated by more than a
year of development and by a different recording session, so the leakage is
limited, but it is not zero, and the cross-patient figures should be read as
very slightly optimistic on that account. Treating them as one unit instead
would remove the concern at the cost of one fewer held-out unit; we report
which fold each falls in with the released per-fold results so that either
convention can be recovered.

\textbf{Siena Scalp EEG Database.} Fourteen adult subjects (20--71 years;
thirteen with focal impaired-awareness seizures, one with focal-to-bilateral
tonic--clonic) recorded at the University of Siena at \SI{512}{\hertz} in a
referential montage with 29 EEG channels plus EKG. Seizure times are given in
per-subject text files relative to registration start.

The two cohorts differ in acquisition hardware, sampling rate, montage and
mains frequency, and also in age group, syndrome mix and clinical context. A
transfer between them is therefore a cohort shift in the full sense rather
than an acquisition shift alone, and the two cannot be separated with these
data. We use ``cohort shift'' throughout for that reason and return to the
distinction in \S\ref{sec:whatshift}.

\textbf{Annotation handling.} Siena annotation files carry free text alongside
timestamps in two instances: an alternative time introduced by ``oppure'', and
separate clinical and electrographic onsets. We parse the first well-formed
\texttt{HH.MM.SS} token in each field, which preserves the clinical onset. One
entry (PN00-3) implies a seizure of \SI{3660}{\second} and is treated as a
transcription error; 43 of 44 annotated seizures are retained. The exclusion
is recorded in the released manifest.

\subsection{Preprocessing}

\textbf{Montage harmonisation.} We derived the 17 bipolar channels of the
CHB-MIT longitudinal montage in both cohorts. For Siena this required
constructing bipolar derivations from the referential recording using the
older T3/T4/T5/T6 nomenclature those files carry. All 41 Siena recordings and
683 of 686 CHB-MIT recordings contain the electrodes required for all 17
derivations. CHB-MIT contains a duplicated \texttt{T8-P8} channel, which we
exclude to avoid weighting one derivation twice.

\begin{figure}[t]
\centering
\includegraphics[width=\textwidth]{figures/F6_spectra.pdf}
\caption{\textbf{The two cohorts handle mains interference in opposite ways,
which sets the usable band.} (a) CHB-MIT already carries a notch: power at
\SI{60}{\hertz} lies about an order of magnitude below the surrounding
spectrum and rises again by \SI{70}{\hertz}, which cannot occur in an
unfiltered $1/f$ spectrum; the stopband spans roughly
\SIrange{50}{70}{\hertz}. (b) Siena is un-notched, with a sharp
\SI{50}{\hertz} peak and a harmonic at \SI{100}{\hertz}. The shaded
\SIrange{1}{40}{\hertz} band lies below both artefacts. Fitting a spectral
model across the CHB-MIT stopband inflates the estimated aperiodic exponent
by approximately \SI{50}{\percent}.}
\label{fig:spectra}
\end{figure}


\textbf{Spectral considerations.} The two cohorts handle mains interference in
opposite ways, and this constrains the usable frequency range
(figure~\ref{fig:spectra}). In CHB-MIT a notch has already been applied: power
at \SI{60}{\hertz} sits roughly an order of magnitude \emph{below} the
surrounding spectrum and rises again by \SI{70}{\hertz}, which cannot occur in
an unfiltered $1/f$ spectrum; the stopband spans approximately
\SIrange{50}{70}{\hertz}. Siena is un-notched and shows a sharp
\SI{50}{\hertz} peak, about fifteen times the adjacent spectrum, with a
harmonic at \SI{100}{\hertz}.

We therefore band-pass filtered both cohorts to \SIrange{1}{40}{\hertz}, which
lies below both artefacts. This constraint is not incidental: fitting a
spectral model across the CHB-MIT stopband inflates the estimated aperiodic
exponent by approximately \SI{50}{\percent} (mean 2.03 over
\SIrange{4}{45}{\hertz} against 3.14 over \SIrange{4}{55}{\hertz} in the same
recordings), an artefact of the filter rather than a property of the signal.

\textbf{Resampling and referencing.} Siena recordings were resampled from
\SI{512}{\hertz} to \SI{256}{\hertz} after filtering. Each channel was
z-scored within its own recording, which removes per-session amplitude scaling
without using information from any other recording.

\subsection{Windowing and labelling}

Recordings were divided into non-overlapping \SI{4}{\second} windows (1024
samples). We use no overlap: overlapping windows drawn from the same seizure
are near-duplicates, and their presence on both sides of a split inflates
apparent performance.

\begin{figure}[t]
\centering
\includegraphics[width=\textwidth]{figures/F9_labelling.pdf}
\caption{\textbf{Window labelling.} A window is positive when at least half
of it falls inside an annotated seizure and negative when none of it does.
Windows with intermediate overlap are discarded, since their label is
genuinely ambiguous and assigning them either way introduces avoidable noise
at the seizure boundary. Windows do not overlap one another.}
\label{fig:labelling}
\end{figure}


A window was labelled positive if at least \SI{50}{\percent} of its duration
fell within an annotated seizure, and negative if no part of it did. Windows
with intermediate overlap were discarded rather than assigned to either class,
since their label is genuinely ambiguous (figure~\ref{fig:labelling}). This
yields \num{1000311} labelled windows across 38 subjects (CHB-MIT
\num{873428}; Siena \num{126883}) with \num{3405} positive windows, a native
positive rate of \SI{0.34}{\percent}.

\subsection{Subsampling for the learned models}

Retaining every window as raw time series would require approximately
\SI{70}{\giga\byte}. We constructed a reduced set comprising every positive
window and, from each recording, 60 randomly selected negative windows:
\num{46245} windows at \SI{7.36}{\percent} positive. All seizure data is
retained; only interictal data is reduced. This changes the evaluation
prevalence, which we treat as a variable to be studied
(section~\ref{sec:prev}) rather than a nuisance to be hidden.

\subsection{Feature families and models}

\textbf{Handcrafted features} (119 per window). Relative band power in five
bands (delta \SIrange{1}{4}{\hertz}, theta \SIrange{4}{8}{\hertz}, alpha
\SIrange{8}{13}{\hertz}, beta \SIrange{13}{30}{\hertz}, gamma
\SIrange{30}{40}{\hertz}) for each of 17 channels (85 features), and the
aperiodic exponent with its goodness of fit from spectral parameterisation
\citep{donoghue2020} over \SIrange{1}{40}{\hertz} for each channel (34
features). Power spectra were computed by Welch's method with a
\SI{1}{\second} Hann window and \SI{50}{\percent} overlap. The fixed aperiodic
model was used; the knee model did not converge on this data, returning
exponents outside any physiologically plausible range. Classification used
logistic regression with balanced class weights on standardised features.

\textbf{Learned representations.} Fourteen published architectures were used
as implemented in braindecode \citep{schirrmeister2017}, without modification,
spanning \num{2402} to \num{1141889} parameters: EEGNet \citep{lawhern2018},
EEGITNet \citep{salami2022}, EEGTCNet \citep{ingolfsson2020}, SincShallowNet
\citep{borra2020}, SCCNet \citep{wei2019}, EEGInceptionERP
\citep{santamaria2020}, ShallowFBCSPNet and Deep4Net
\citep{schirrmeister2017}, EEGNeX \citep{chen2024}, ATCNet
\citep{altaheri2022}, TIDNet \citep{kostas2020}, EEGSimpleConv
\citep{ouahidi2023}, EEGConformer \citep{song2022} and SPARCNet
\citep{jing2023}. Using an existing library rather than reimplementing removes
one source of divergence from the original reports.

Training used AdamW (learning rate \num{e-3}, weight decay \num{e-4}), batch
size 256, class-weighted cross-entropy, mixed precision, and early stopping on
validation AUROC with patience 4 over at most 20 epochs, with
\SI{10}{\percent} of each training split held out for validation. No
hyperparameter search was performed and identical settings were used for every
architecture; the objective is comparison under a common protocol, not
per-model optimisation.

\subsection{Generalisation settings}

\begin{figure}[t]
\centering
\includegraphics[width=\textwidth]{figures/F1_pipeline.pdf}
\caption{\textbf{Study design.} Two cohorts differing in age group and
clinical context as well as in acquisition hardware, sampling rate, montage
and mains frequency are harmonised to a common 17-channel bipolar montage, band-passed to
\SIrange{1}{40}{\hertz}, resampled to \SI{256}{\hertz} and z-scored within
each recording. Windows are \SI{4}{\second} with no overlap; windows
overlapping a seizure by less than half are discarded rather than labelled.
The same windows, splits and preprocessing are used for every model, so
differences between the three settings are attributable to the setting
rather than to the pipeline.}
\label{fig:pipeline}
\end{figure}


Three nested settings were evaluated (figure~\ref{fig:pipeline}). \textbf{S1},
within-patient: the first \SI{70}{\percent} of one subject's windows in
temporal order for training, the remainder for testing. \textbf{S2},
cross-patient: grouped five-fold cross-validation with folds formed over
CHB-MIT subjects, so no subject contributes to both sides of a split.
\textbf{S3}, cross-cohort: all of one cohort for training and all of the other
for testing, in both directions. The comparison of interest is the S1$\to$S2
drop against the S2$\to$S3 drop: the first is the cost of a new patient, the
second the additional cost of a new centre. Only one of the two S3
directions nests cleanly inside S1 and S2. Those two settings are both
trained on CHB-MIT, so CHB-MIT $\rightarrow$ Siena continues the same
sequence from the same source cohort, whereas Siena $\rightarrow$ CHB-MIT
changes the training cohort as well and is not a further step in that
sequence. We therefore take CHB-MIT $\rightarrow$ Siena as the primary
comparison and report Siena $\rightarrow$ CHB-MIT separately, as a
replication in the reverse direction.

For the learned models, S1 was run on every CHB-MIT subject with at least 30
positive windows on the training side and 8 on the test side. This admits 8
of the 24 recording sets. The restriction is not only one of sample size: in several
subjects every annotated seizure falls in the final \SI{30}{\percent} of the
recording, so a temporal within-subject split leaves no positive training
examples at all, and those subjects cannot enter S1 under any threshold.

\subsection{Prevalence sweep}\label{sec:prev}

To quantify protocol dependence we held the trained model fixed and varied
only the class balance of the evaluation set. All positives were retained and
negatives were drawn without replacement to achieve target positive rates of
50, 25, 10, 7.36, 5 and \SI{1}{\percent}, together with the native rate of
\SI{0.34}{\percent}. The upper end of this range corresponds to the balanced
evaluation common in the literature; the lower end to continuous monitoring.

\subsection{Seed variance and calibration}

Because the sweep reports differences between architectures of the order of
0.02--0.05 AUROC, we characterised the variation attributable to
initialisation alone. Four architectures spanning 2.4$\times$10$^3$ to
1.1$\times$10$^6$ parameters (EEGNet, ShallowFBCSPNet, ATCNet, SPARCNet)
were rerun in every setting with five random seeds, holding the data,
splits and hyperparameters fixed: 300 runs. This allows the spread between
architectures to be set directly against the spread produced by
initialisation alone, which is the quantity the paper's central claim turns
on.
Calibration was quantified per run by the Brier score and by expected
calibration error over fifteen equal-width bins.

\subsection{Negative controls}

\textbf{Label permutation.} Labels were permuted within each subject,
preserving per-subject class balance, and the full pipeline re-run, three
times.

\textbf{Recording context.} In CHB-MIT all positive windows originate from
seizure-containing recordings. If those recordings carry a session-level
signature---electrode adjustment, arousal, staff activity---a classifier could
score well without representing ictal activity. We therefore trained a
classifier to separate \emph{non-seizure} windows of seizure-containing
recordings from \emph{non-seizure} windows of seizure-free recordings, using
folds grouped over subjects.

\subsection{Feature-space structure}

To characterise what the features encode, we trained logistic regression to
recover three targets from the same windows under identical five-fold
cross-validation: cohort identity (binary), subject identity (one-vs-rest,
averaged over subjects with at least 50 windows), and the seizure label.
Silhouette scores were computed on the first two principal components.

\subsection{Analyses from the saved predictions}

The score of every test window was retained, with its subject and its
position in the recording, so that the following could be computed without
retraining.

\emph{Matched prevalence.} The learned models were trained on a subsample
of the windows (\SI{7.36}{\percent} positive) while the handcrafted features
were evaluated on everything (\SI{0.34}{\percent}), which is not a
like-for-like comparison. We addressed this in two ways. Each model's own
test windows were resampled to a range of positive rates with the trained
model held fixed; and, separately, every window of both cohorts was
extracted as raw time series and the trained models were applied to all of
it, so that the two representations could be compared on exactly the same
\num{1000311} windows at the prevalence of continuous recording. The second
required no retraining: the weights saved during the sweep were reused.

\emph{Subject-level bootstrap.} Windows are not independent observations:
tens of thousands of them come from a few dozen people, and resampling
windows would give intervals far narrower than the estimate deserves. We
resampled subjects with replacement instead, moving all of a subject's
windows together, over 2000 replicates.

\emph{Event-level evaluation.} Contiguous positive predictions were merged
into events. An annotated seizure counted as detected if any predicted event
overlapped it; every other predicted event counted towards the false-alarm
rate. Sensitivity is reported at fixed false alarms per hour. Because the
windows are a subsample of continuous recording, the denominator is
effective rather than wall-clock time, so the false-alarm rate is optimistic
and is comparable between the systems evaluated here rather than against
figures obtained on full recordings.

\emph{Reliability.} Predicted probabilities were binned into fifteen
equal-width intervals and the observed frequency of the positive class
computed within each.

\subsection{Metrics and statistics}

We report AUROC and AUPRC. At a native positive rate of \SI{0.34}{\percent},
accuracy and specificity are uninformative, and AUROC alone is insufficient
because it is invariant to prevalence while the precision achievable at a
given sensitivity is not; AUPRC is therefore reported alongside the prevalence
against which it should be read. The S1$\to$S2 and S2$\to$S3 drops were
compared across architectures three ways: by resampling architectures with
replacement (20\,000 draws) to obtain intervals on each drop and on their
ratio; by a paired permutation test that randomises the sign of each
architecture's difference; and by the Wilcoxon signed-rank and paired
$t$-tests. These are computed on the nested direction; the reverse direction
is reported the same way but separately. The ratio is reported alongside the difference because it is the
better estimated of the two when both quantities are small.

\subsection{Implementation and availability}

Analyses used Python 3.12 with MNE-Python \citep{gramfort2013}, NeuroDSP,
specparam \citep{donoghue2020}, scikit-learn, PyTorch and braindecode, with
versions pinned in the released environment file. Learned models were trained
on a single NVIDIA T4. Code, the window-level manifest and all per-fold
results are available at \url{https://github.com/<user>/eeg-cross-cohort-seizure-detection}\\
and archived at \url{https://doi.org/<Zenodo DOI on release>}.

%% =====================================================================
\section{Results}
%% =====================================================================

\subsection{Negative controls}

\begin{figure}[t]
\centering
\includegraphics[width=0.52\textwidth]{figures/F5_controls.pdf}
\caption{\textbf{Negative controls.} Detection performance with the actual
labels, against two controls that should sit at chance and do. Permuting
labels within subjects destroys the label--signal relationship; the
recording-context control asks whether non-seizure windows of
seizure-containing recordings can be separated from non-seizure windows of
seizure-free recordings, which would indicate that apparent detection
performance partly reflects session-level differences rather than ictal
activity. Error bars are standard deviations across folds and repeats.}
\label{fig:controls}
\end{figure}


Label permutation gave a mean AUROC of 0.511 across three repeats and both S2
and S3 settings (range 0.480--0.541), with AUPRC 0.0013--0.0046 against a base
rate of 0.0034. Performance collapses to chance when the label--signal
relationship is destroyed.

The recording-context test gave a mean AUROC of 0.522 across five grouped
folds (range 0.476--0.572) for separating non-seizure windows of
seizure-containing recordings ($n=\num{165819}$) from non-seizure windows of
seizure-free recordings ($n=\num{704791}$). Session-level context is not
linearly separable in these features, and the detection performance reported
below is therefore unlikely to be an artefact of recording context
(figure~\ref{fig:controls}).

\subsection{Handcrafted features under transfer}\label{sec:hand}

Table~\ref{tab:linear} gives leave-one-subject-out and cross-cohort
performance for the handcrafted features on the full continuous data.
Combining the two families outperformed either alone in every setting.

\begin{table}[t]
\centering\small
\caption{Handcrafted features, continuous data at native prevalence (AUROC).
Leave-one-subject-out values are means over held-out subjects.}
\label{tab:linear}
\begin{tabular}{lccc}
\toprule
Setting & band power & aperiodic & both \\
\midrule
Leave-one-out, CHB-MIT ($n=24$ recording sets) & 0.613 & 0.636 & 0.718 \\
Leave-one-subject-out, Siena ($n=12$)   & 0.706 & 0.665 & 0.747 \\
CHB-MIT $\rightarrow$ Siena             & 0.638 & 0.677 & 0.687 \\
Siena $\rightarrow$ CHB-MIT             & 0.612 & 0.617 & 0.680 \\
\bottomrule
\end{tabular}
\end{table}

The comparison the study turns on appears directly in these numbers. With
band-power features, 13 of 24 CHB-MIT leave-one-out folds scored
\emph{worse} than the cross-cohort transfer, and 11 of 24 in the reverse
direction; with the aperiodic features, 16 of 24 and 10 of 24; with the
combined set, 8 of 24 and 7 of 24. In a substantial fraction of cases, testing
on a new patient from the same centre is harder than testing on a different
centre entirely.

\subsection{Learned representations under transfer}\label{sec:arch}

\begin{figure}[t]
\centering
\includegraphics[width=\textwidth]{figures/F2_architectures.pdf}
\caption{\textbf{Generalisation across fourteen architectures.} (a) Each
line is one architecture across the three settings; the heavy line is their
mean and the horizontal marks give the handcrafted baseline at S2 and S3.
The separation between learned and handcrafted representations under cohort
transfer exceeds the spread among the learned representations themselves.
(b) The two costs per architecture. The S1$\to$S2 drop, the cost of a new
patient, exceeds the S2$\to$S3 drop, the additional cost of a new centre, in
11 of 14 architectures. Two architectures return S1 below S2, which cannot
reflect a genuine ordering and instead reflects the instability of the S1
estimate.}
\label{fig:arch}
\end{figure}


Table~\ref{tab:dl} gives the fourteen architectures across the three settings,
and figure~\ref{fig:arch} shows the same data together with the handcrafted
reference.

\begin{table}[t]
\centering\small
\caption{Fourteen architectures across the three settings (mean AUROC).
Architectures are ordered by parameter count. The final two columns give the
cost of each boundary.}
\label{tab:dl}
\begin{tabular}{lrccccc}
\toprule
Architecture & Params & S1 & S2 & S3 & S1$\to$S2 & S2$\to$S3 \\
\midrule
EEGNet           & \num{2402}    & 0.937 & 0.850 & 0.845 &  0.087 &  0.005 \\
EEGITNet         & \num{3670}    & 0.932 & 0.878 & 0.842 &  0.054 &  0.036 \\
EEGTCNet         & \num{4198}    & 0.823 & 0.875 & 0.865 & $-$0.052 & 0.010 \\
SincShallowNet   & \num{11714}   & 0.822 & 0.881 & 0.831 & $-$0.060 & 0.050 \\
SCCNet           & \num{12882}   & 0.957 & 0.856 & 0.810 &  0.102 &  0.046 \\
EEGInceptionERP  & \num{27342}   & 0.953 & 0.835 & 0.785 &  0.118 &  0.050 \\
ShallowFBCSPNet  & \num{33282}   & 0.961 & 0.866 & 0.811 &  0.096 &  0.054 \\
EEGNeX           & \num{55650}   & 0.978 & 0.851 & 0.830 &  0.127 &  0.020 \\
ATCNet           & \num{113242}  & 0.953 & 0.888 & 0.838 &  0.066 &  0.050 \\
TIDNet           & \num{200274}  & 0.870 & 0.852 & 0.834 &  0.018 &  0.018 \\
Deep4Net         & \num{277352}  & 0.929 & 0.816 & 0.818 &  0.113 & $-$0.002 \\
EEGSimpleConv    & \num{724922}  & 0.944 & 0.873 & 0.828 &  0.071 &  0.045 \\
EEGConformer     & \num{791746}  & 0.953 & 0.868 & 0.849 &  0.085 &  0.019 \\
SPARCNet         & \num{1141889} & 0.924 & 0.860 & 0.848 &  0.064 &  0.012 \\
\midrule
\textbf{mean}    &               & \textbf{0.924} & \textbf{0.861} & \textbf{0.831} & \textbf{0.064} & \textbf{0.029} \\
\addlinespace
handcrafted (both) & 119 features & --- & 0.709 & 0.684 & --- & --- \\
\bottomrule
\end{tabular}
\end{table}

Three observations follow.

\textbf{Representation separates the two families; architecture does not.}
Under cross-cohort transfer the fourteen learned representations averaged
0.831 against 0.684 for the handcrafted features, a difference of 0.147. The
spread among the learned representations themselves was 0.785--0.865, so
moving from feature engineering to a learned representation mattered
substantially more than the choice among learned representations.

\textbf{The cohort boundary costs less than the patient boundary, by an
amount this design does not pin down.} In the nested comparison --- S1 and
S2 both trained on CHB-MIT, then CHB-MIT $\rightarrow$ Siena --- the mean
S1$\to$S2 drop was 0.063 (bootstrap \SI{95}{\percent} CI 0.032 to 0.090) and
the mean S2$\to$S3 drop 0.039 (0.030 to 0.049), a ratio of 0.62 with an
interval of 0.37 to 1.35. The direction was consistent in 10 of 14
architectures, but the interval includes unity and the paired difference of
0.024 ($-$0.010 to 0.055) does not reach conventional significance (paired
permutation $p=0.197$; Wilcoxon $p=0.153$).

The reverse direction, Siena $\rightarrow$ CHB-MIT, gives a smaller cost
still: an S2$\to$S3 drop of 0.020 (0.002 to 0.038), a ratio of 0.32 (0.02 to
0.72), consistent in 12 of 14 architectures. This is a replication rather
than a further step in the sequence, since the training cohort changes, but
it points the same way and more strongly.

Pooling the two directions, as an earlier version of this analysis did,
gives a ratio of 0.47 (0.27 to 0.98) and an interval that excludes unity.
We report the nested figure as primary because pooling mixes a nested step
with one that is not, and because the pooled estimate is the more favourable
of the two --- which is precisely the kind of protocol choice this paper
argues should be made explicit. The conclusion the evidence supports is that
the cohort boundary is consistently the cheaper of the two, at somewhere
between a third and roughly all of the patient boundary's cost, and that
this design does not resolve it more finely than that.

Two architectures, EEGTCNet and SincShallowNet, returned S1 below S2, which
cannot reflect a genuine ordering and instead reflects the instability of
the S1 estimate, computed on four subjects and depressed in both cases by a
single subject (chb13: 0.544 and 0.442 respectively). Excluding them as a
sensitivity analysis widens the S1$\to$S2 drop to 0.083 and leaves the
S2$\to$S3 drop unchanged; we report this as sensitivity rather than as the
primary result, since the exclusion was identified after inspecting the
values.

\textbf{Absolute performance remains modest at native prevalence.} The
learned representations were evaluated at \SI{7.36}{\percent} positive; the
handcrafted features at \SI{0.34}{\percent}. The next section separates the
contribution of that difference.

\subsection{Comparison at matched prevalence}

Two comparisons place the representations on the same footing.

Resampling each learned model's own test windows to a common positive rate
removes the difference in class balance. At \SI{7.36}{\percent}, the learned
representations reached AUROC 0.868 in S2 and 0.853 in S3, against 0.709 and
0.687 for the handcrafted features on the matched-window subset
(\S\ref{sec:hand}); a separation of roughly 0.16 attributable to the
representation rather than to the prevalence (figure~\ref{fig:matched}).

\begin{figure}[t]
\centering
\includegraphics[width=\textwidth]{figures/F14_native_comparison.pdf}
\caption{\textbf{Both representations on the full continuous set.} The
trained models are applied without retraining to every window of the
relevant test set, and the evaluation prevalence is then varied from
balanced down to the native \SI{0.34}{\percent} (dashed line). (a) AUROC is
flat across the range and separates the two representations by about 0.16 at
every point. (b) AUPRC falls with the base rate for both, but the ratio
between them widens, because a lift of 42.5$\times$ and one of 2.8$\times$
diverge in absolute terms as the base rate falls. The comparison is
like-for-like: the same recordings, the same windows and the same
prevalence.}
\label{fig:nativecmp}
\end{figure}


Applying the trained models to every window of both cohorts settles it
directly. On the same \num{1000311} windows and at the same
\SI{0.34}{\percent} prevalence, the four architectures reached a mean
cross-cohort AUROC of 0.843 against 0.686 for the handcrafted features
(table~\ref{tab:native}, figure~\ref{fig:nativecmp}). The gap is 0.157, in the same direction and of
the same size as the matched-subset estimate. Expressed as lift over chance,
which is what determines whether a detector is usable at this base rate, the
difference is larger: 42.5$\times$ against 2.8$\times$.

\begin{table}[t]
\centering\small
\caption{Both representations on continuous data at native prevalence, with
no retraining: the models are those from the sweep applied to every window
of the relevant test set. Learned figures are means over four architectures
and five seeds. Lift is AUPRC divided by the positive-class rate of the set
in question, and is the quantity that determines whether a detector is
usable at a base rate of this order.}
\label{tab:native}
\begin{tabular}{lcccc}
\toprule
& AUROC & AUPRC & lift & recording \\
\midrule
\multicolumn{5}{l}{\emph{learned representations}} \\
\quad S1 within-patient          & 0.961 & 0.491 & 118.7$\times$ & 37 h \\
\quad S3 cross-cohort            & 0.843 & 0.155 &  42.5$\times$ & 556 h \\
\qquad CHB-MIT $\rightarrow$ Siena & 0.828 & 0.122 &  26.3$\times$ & 141 h \\
\qquad Siena $\rightarrow$ CHB-MIT & 0.859 & 0.189 &  58.7$\times$ & 970 h \\
\addlinespace
\multicolumn{5}{l}{\emph{handcrafted features, same windows}} \\
\quad S2 cross-patient           & 0.693 & 0.011 &   3.6$\times$ & --- \\
\quad S3 cross-cohort            & 0.686 & 0.011 &   2.8$\times$ & 556 h \\
\bottomrule
\end{tabular}
\end{table}

\begin{figure}[t]
\centering
\includegraphics[width=0.62\textwidth]{figures/F15_event_native.pdf}
\caption{\textbf{Event-level operating curves on continuous recording.}
Sensitivity against false alarms per hour, with the denominator real
recording time rather than a subsample, so the two representations and the
two settings can be read against each other and against figures reported
elsewhere. Dotted lines mark 0.1 and 1 false alarms per hour. The
within-patient curve approaches a usable operating point; the cross-cohort
curve does not, and the handcrafted curve is lower again by roughly a factor
of twenty at one false alarm per hour.}
\label{fig:eventnative}
\end{figure}


Scored at the event level on the same continuous data, where the
false-alarm denominator is real recording time, the separation persists
(table~\ref{tab:eventnative}, figure~\ref{fig:eventnative}). At one false alarm per hour the learned
representations detected \SI{22.1}{\percent} of cross-cohort seizures
against \SI{1.1}{\percent} for the handcrafted features; at two per hour,
\SI{49.8}{\percent} against \SI{17.0}{\percent}. Median detection delay was
16 s against 28 s. Within patients the same models reached
\SI{77.8}{\percent} at one false alarm per hour, with an 8 s delay.

\begin{table}[t]
\centering\small
\caption{Event-level sensitivity on continuous data. Contiguous positive
predictions are merged into events and a seizure counts as detected if any
predicted event overlaps it. Unlike table~\ref{tab:event}, the denominator
here is real recording time, so these rates are comparable between the two
representations and with figures reported elsewhere on full recordings.}
\label{tab:eventnative}
\begin{tabular}{lccc}
\toprule
False alarms per hour & learned S1 & learned S3 & handcrafted S3 \\
\midrule
0.1 & 0.166 & 0.067 & 0.011 \\
0.5 & 0.504 & 0.091 & 0.011 \\
1.0 & 0.778 & 0.221 & 0.011 \\
2.0 & 0.959 & 0.498 & 0.170 \\
5.0 & 0.985 & 0.646 & 0.290 \\
\addlinespace
median delay & 8 s & 16 s & 28 s \\
\bottomrule
\end{tabular}
\end{table}

The full-data evaluation also reproduces the architecture result on
independent ground. Transferring across cohorts, the four architectures
spanned 0.835 to 0.855, a range of 0.021 against a seed standard deviation
of 0.020: a ratio of 1.05, against 1.09 on the subsampled set. A different
dataset size, a different prevalence and a different evaluation return the
same conclusion about architecture.

AUROC was, as expected, largely unmoved by the resampling: 0.869 at
\SI{50}{\percent} against 0.868 at \SI{7.36}{\percent} in S2. AUPRC moved
substantially over the same range, from 0.895 to 0.658, while the lift over
chance rose from 1.8$\times$ to 8.9$\times$. The two metrics move in opposite
directions as prevalence falls because AUPRC tracks the base rate while lift
does not, which is the reason a raw AUPRC cannot be read without it.

\subsection{Event-level performance}

Scoring the same predictions as a clinical system would be scored separates
the settings far more sharply than window-level ranking does
(table~\ref{tab:event}, figure~\ref{fig:event}). Within patients,
event-level sensitivity reached 0.832 at one false alarm per hour. Across
patients it was 0.413 and across cohorts 0.419 --- again close to one
another, reproducing at the event level the pattern found at the window
level, under a metric that makes no use of the ranking's tail.

\begin{table}[t]
\centering\small
\caption{Event-level sensitivity at fixed false-alarm rates. Contiguous
positive predictions are merged into events; a seizure counts as detected if
any predicted event overlaps it. Because the evaluation windows are a
subsample of continuous recording, these false-alarm rates are optimistic.}
\label{tab:event}
\begin{tabular}{lccc}
\toprule
False alarms per hour & S1 within-patient & S2 cross-patient & S3 cross-cohort \\
\midrule
0.1 & 0.832 & 0.119 & 0.182 \\
0.5 & 0.832 & 0.325 & 0.299 \\
1.0 & 0.832 & 0.413 & 0.419 \\
2.0 & 0.894 & 0.478 & 0.532 \\
\bottomrule
\end{tabular}
\end{table}

Two things follow. A window-level AUROC of about 0.85 corresponds to
detecting fewer than half of seizures at one false alarm per hour, so the
choice of framing changes the apparent maturity of a system by more than any
difference this study finds between architectures or between
representations. And the cross-patient and cross-cohort figures agree to
within 0.006 at this operating point, which is a stronger form of the
paper's central result than the window-level comparison, since it holds
under a metric that played no part in selecting anything.

\subsection{Prevalence dependence}

\begin{figure}[t]
\centering
\includegraphics[width=\textwidth]{figures/F3_prevalence.pdf}
\caption{\textbf{Evaluation prevalence changes the reported figure without
changing the model.} The trained model is held fixed and only the class
balance of the evaluation set is varied. (a) AUROC moves by 0.033 across the
whole range. (b) AUPRC moves by a factor of 165 on the same runs, tracking
the base rate shown dotted. The shaded region marks the balanced evaluation
common in the literature; the dashed line marks the native
\SI{0.34}{\percent} of continuous recording. A figure reported in one regime
and a figure reported in the other are not comparable on AUPRC, and AUROC
does not reveal the difference.}
\label{fig:prev}
\end{figure}


With the trained model held fixed and only the evaluation class balance
varied, AUROC moved by 0.033 across the full range while AUPRC moved by a
factor of 165 (0.810 at \SI{50}{\percent} prevalence, 0.005 at the native
\SI{0.34}{\percent}); see figure~\ref{fig:prev}. Precision lift over chance
remained between 1.6$\times$ and 3.9$\times$ throughout, so the apparent
collapse reflects the changing base rate rather than a change in
discriminative ability.

The practical consequence is that a system reported at balanced prevalence and
a system reported at native prevalence cannot be compared on AUPRC, and AUROC
does not reveal the difference. The range swept here spans the protocols in
current use.

\subsection{What the features encode}

\begin{figure}[t]
\centering
\includegraphics[width=\textwidth]{figures/F8_embedding.pdf}
\caption{\textbf{What the features encode.} (a--c) The same principal
component projection coloured by cohort, by subject and by label; the two
components shown carry a minority of the variance, so absence of visible
separation is not absence of signal. (d) Each target recovered from the full
feature set by logistic regression under identical five-fold
cross-validation. Cohort identity and subject identity are similarly
recoverable, and both exceed the seizure label. In this representation the
cohort boundary is not a stronger source of structure than the subject
boundary, which is the mechanism behind the small S2$\to$S3 cost. The
subject probe is one-vs-rest across 38 subjects while the cohort probe is
binary, so the comparison is indicative rather than definitive.}
\label{fig:embed}
\end{figure}


Measured identically---five-fold logistic regression on the same
windows---cohort identity was recoverable at AUROC 0.92, subject identity at
0.95, and the seizure label itself at 0.81 (figure~\ref{fig:embed}). Both
forms of identity are more recoverable than the label the model is meant to
predict, and the two are close to each other.

The per-feature magnitude of cohort shift, measured by the Kolmogorov--Smirnov
statistic between cohorts, was uncorrelated with that feature's importance for
seizure detection ($r=-0.096$). Shift was concentrated in the beta band (mean
KS 0.249) and in central and midline derivations, while the strongest
discriminative weights sat elsewhere. Mean KS was similar for the two feature
families (band power 0.185, aperiodic 0.181).

\begin{figure}[t]
\centering
\includegraphics[width=\textwidth]{figures/F10_matched_prevalence.pdf}
\caption{\textbf{Learned and handcrafted representations compared at the same
class balance.} Each model's own test windows are resampled to a range of
positive rates with the trained model held fixed. (a) AUROC is essentially
unmoved. (b) AUPRC falls with the base rate while the lift over chance
rises, which is why AUPRC cannot be read without the prevalence it was
computed at. The separation from the handcrafted baseline persists at every
matched prevalence.}
\label{fig:matched}
\end{figure}

\begin{figure}[t]
\centering
\includegraphics[width=0.62\textwidth]{figures/F11_event_level.pdf}
\caption{\textbf{Event-level sensitivity against false alarms per hour.}
Contiguous positive predictions are merged into events. The within-patient
curve sits far above the cross-patient and cross-cohort curves, which are
close to one another. The false-alarm rates are optimistic because the
evaluation windows are a subsample of continuous recording.}
\label{fig:event}
\end{figure}

\begin{figure}[t]
\centering
\includegraphics[width=\textwidth]{figures/F12_reliability.pdf}
\caption{\textbf{Reliability diagrams by setting.} Observed frequency of the
positive class against predicted probability, with the diagonal marking
perfect calibration. The curves fall below the diagonal in all three
settings: the models are over-confident throughout, not only under transfer.}
\label{fig:reliab}
\end{figure}


\subsection{Per-subject variation}

\begin{figure}[t]
\centering
\includegraphics[width=\textwidth]{figures/F4_per_subject.pdf}
\caption{\textbf{Per-subject performance, which the cohort mean conceals.}
(a) Leave-one-subject-out AUROC for each held-out CHB-MIT subject, ordered
by the combined feature set. The shaded region lies below chance. (b) The
same values as distributions. Individual subjects fall well below chance
--- chb14 reaches 0.141 with band power, indicating systematically inverted
ranking rather than absent signal. Combining feature families narrows the
spread but does not remove the failures.}
\label{fig:persub}
\end{figure}


Leave-one-subject-out performance varied widely and included results below
chance that cohort means conceal (figure~\ref{fig:persub}). With band power, 8
of the 24 recording sets fell below 0.5 and chb14 reached 0.141, indicating
systematically inverted ranking rather than absent signal. Combining feature
families reduced this to 2 of 24 and narrowed the spread (SD 0.124 against
0.192) but did not eliminate it: chb13 and chb15 remained below 0.5 in all
three feature sets. The same pattern appears among the learned models, where
chb13 gave 0.442 with SincShallowNet and 0.917 with SCCNet.

\subsection{Variation attributable to initialisation}

Four architectures rerun with five seeds each, with everything else held
fixed, answer directly how much of the difference between them is real
(table~\ref{tab:seed}). Averaged over configurations, initialisation alone
moved AUROC by 0.030 in S1, 0.017 in S2 and 0.019 in S3.

Set against that, the spread between the four architectures was 0.012 in S2
and 0.021 in S3, giving ratios of architecture spread to seed noise of 0.66
and 1.09. In the cross-patient setting the four architectures are separated
by less than the variation a single one of them shows across seeds; in the
cross-cohort setting, by about the same amount.

The consequence is visible in the ranking. Which architecture performed best
cross-cohort changed with the seed: SPARCNet at seed 1234, ATCNet at 3456
and 7890, EEGNet at 5678 and 9012. Three of the four won at least once
across five runs that differed in nothing but initialisation. Cross-patient,
ATCNet led at four seeds of five and ShallowFBCSPNet at the fifth.

Only within patients does architecture separate the models reliably (spread
0.071, ratio 2.39), and that is the setting furthest from deployment. The
within-patient instability is concentrated in one architecture rather than
being a property of the setting: EEGNet gave a mean per-subject seed
standard deviation of 0.072 in S1, with subject chb12 ranging from 0.451 to
0.837 across five seeds, while ShallowFBCSPNet gave 0.009 in the same
setting --- lower than its own cross-patient and cross-cohort figures.
Reproducibility is therefore a property of the architecture-and-setting
pair, not of the setting alone.

Two consequences follow. Differences between architectures of the order of
0.02--0.03 AUROC in S2 and S3 are of the same magnitude as the seed
variation and should not be read as architectural. And the S1 column of
table~\ref{tab:dl}, computed from a single seed on eight subjects, is the
least reliable part of the design; the two architectures returning S1 below
S2 are best understood in that light rather than as a property of those
models.

\begin{table}[t]
\centering\small
\caption{Variation attributable to initialisation set against variation
attributable to architecture. Four architectures (EEGNet, ShallowFBCSPNet,
ATCNet, SPARCNet), five seeds each, in every setting, with data, splits and
hyperparameters held fixed: 300 runs. A ratio at or below one means the
architectures are separated by no more than the noise from initialisation.}
\label{tab:seed}
\begin{tabular}{lcccc}
\toprule
Setting & mean seed SD & architecture spread & ratio & best architecture \\
\midrule
S1 within-patient & 0.030 & 0.071 & 2.39 & SPARCNet (all seeds) \\
S2 cross-patient  & 0.017 & 0.012 & 0.66 & ATCNet (4 of 5 seeds) \\
S3 cross-cohort   & 0.019 & 0.021 & 1.09 & changes with seed \\
\bottomrule
\end{tabular}
\end{table}

\subsection{Calibration}

Expected calibration error was highest within patients (mean 0.283, range
0.101--0.461), lower across patients (0.166, 0.119--0.233) and lowest across
cohorts (0.127, 0.037--0.232); the Brier score followed the same ordering
(0.151, 0.107, 0.091). Calibration therefore did not deteriorate
under transfer; it was poorest in the setting where ranking performance was
best. The models are miscalibrated from the outset rather than becoming so
when the domain changes, and the class-weighted training objective, which
optimises ranking rather than probability, is the likely reason. Any
deployment using a fixed probability threshold would require recalibration
regardless of whether the data came from a new patient or a new centre.


Predicted probabilities were over-confident in every setting: the mean
predicted probability was 0.291 against an observed positive rate of 0.139
within patients, 0.170 against 0.065 across patients, and 0.364 against
0.074 across cohorts (figure~\ref{fig:reliab}). The pattern is informative
rather than merely unfortunate. Miscalibration of this size was already
present in the easiest setting, so it is not produced by transfer; the
class-weighted objective optimises ranking rather than probability, and the
resulting bias is a property of the training procedure that every one of
these architectures inherits. It also follows that the transfer results
reported above are not confounded by calibration, since ranking metrics are
invariant to a monotone distortion of the scores.

\subsection{Normalisation}

Per-subject standardisation of the features improved transfer in one direction
and degraded it in the other (CHB-MIT $\rightarrow$ Siena 0.687 to 0.745;
Siena $\rightarrow$ CHB-MIT 0.680 to 0.632). Per-subject rank normalisation
behaved similarly (0.739 and 0.666). We report this as an observation rather
than as a recommendation.

%% =====================================================================
\subsection{Uncertainty of the estimates}

Bootstrapping over subjects rather than windows places the estimates on the
scale of the population they generalise to (table~\ref{tab:ci}). Interval
width tracked the number of subjects in the held-out set, as it should: the
mean \SI{95}{\percent} width was 0.107 AUROC across cohorts, where the test
set holds fourteen subjects or twenty-four recording sets, and 0.208 across
patients, where
each fold holds four or five. The cross-cohort estimates, which carry this
study's principal comparison, are therefore the better resolved of the two.

\begin{table}[t]
\centering\small
\caption{Subject-level bootstrap, 2000 replicates over 300 runs. Resampling
windows rather than subjects would treat correlated windows from the same
person as independent and give substantially narrower intervals.}
\label{tab:ci}
\begin{tabular}{lccc}
\toprule
Setting & mean AUROC & mean 95\% CI width & subjects in test set \\
\midrule
S2 cross-patient & 0.870 & 0.208 & 4--5 \\
S3 cross-cohort  & 0.847 & 0.107 & 14 or 24 \\
\bottomrule
\end{tabular}
\end{table}

This bears directly on the central comparison. The mean S1$\to$S2 and
S2$\to$S3 drops of 0.063 and 0.030 are differences between quantities whose
own intervals span more than 0.2, which is consistent with the difference
between them not reaching conventional significance even where the interval
on their ratio excludes unity (\S\ref{sec:arch}), and is the reason we rest
the argument on the effect
size and the consistency of its direction rather than on a test.

\section{Discussion}
%% =====================================================================

\subsection{The dataset boundary is not the expensive one}

Under a single protocol, crossing a cohort boundary cost less than crossing
a patient boundary in the great majority of architectures and in both
directions, and in a substantial fraction of individual folds it cost less
than nothing --- a new patient from the same centre was harder than the
other cohort entirely. How much less is not settled here: the nested
comparison puts the ratio at 0.62 with an interval running from 0.37 to
1.35, and the reverse direction at 0.32. What the evidence supports is the
ordering and its consistency, not a particular multiple.

The mechanism is visible in the feature space. Cohort identity and subject
identity were similarly recoverable (0.92 and 0.95), so in this representation
the two boundaries are alike, and there is no reason to expect one to be
markedly more expensive than the other. This has a direct implication for how
domain adaptation targets are defined. If a method aligns a source dataset to
a target dataset but the residual error is inter-subject, the alignment is
operating on the wrong axis, and the reported improvement would also have been
obtained by pooling more subjects. The observation is consistent with the
survey suggestion that single-source domain adaptation may be theoretically
misplaced for cross-subject EEG \citep{crosssubjsurvey2026}.

A second observation qualifies the first. Cohort shift is real and measurable
here---cohort identity is recoverable at 0.92---but its per-feature magnitude
is uncorrelated with discriminative importance ($r=-0.096$). The features that
move most between the two cohorts are largely not the features that carry the
seizure signal. The shift is therefore present but, in this task and this
band, largely benign. That is a stronger and more useful statement than the
absence of shift would have been.

\subsection{What kind of shift this is}\label{sec:whatshift}

The two cohorts differ in more than acquisition. CHB-MIT is paediatric and
recorded during medication taper in a pre-surgical unit; Siena is adult, and
its syndrome mix is not the same. Sampling rate, montage, amplifier and mains
frequency differ as well. This design cannot separate those factors, and we
do not claim to: what is measured is the cost of moving between two cohorts
that differ in all of these ways at once, which is also the situation a
deployed system faces.

This matters for how the result should be read. That crossing this boundary costs
less than crossing a patient boundary is a statement about cohort transfer
as clinically encountered, not a demonstration that amplifiers and montages
are unimportant. Acquisition differences alone might
cost less, with the population difference carrying most of what we measure;
the two might also partly cancel. Separating them would require cohorts
matched on population but differing in acquisition, or the reverse, and no
such pair is publicly available for scalp EEG at this recording length.

Two observations bear on it indirectly. The shift-importance correlation
above indicates that whatever differs between the cohorts is largely
orthogonal to the seizure signal in the 1--40 Hz band. And the two transfer
directions, which differ by a factor of fourteen in training-set size, gave
similar results, which argues against the measured cost being dominated by
any quantity that scales with data volume.

\subsection{Representation, not architecture}

The gap between handcrafted and learned representations under cohort transfer
(0.147 on the matched subset, 0.157 on the full continuous set) was several
times the spread among fourteen learned architectures spanning three orders
of magnitude in parameter count (0.081). A
\num{2402}-parameter model from 2018 \citep{lawhern2018} transferred as well
as a \num{1141889}-parameter model designed for this exact clinical task
\citep{jing2023}. On this evidence architectural capacity does not appear to
be what limits cross-cohort transfer in this task, though a wider sweep of
capacities and training budgets would be needed to exclude it.

We note that the opposite result has been reported in consumer-grade EEG
emotion recognition, where handcrafted features outperformed end-to-end
learning \citep{consumereeg2026}. The regimes differ in channel count, signal
quality and the amount of labelled data, and the two results together suggest
that the advantage of learned representations is regime-dependent rather than
general. Establishing where the crossover lies would require sweeping those
variables directly.

\subsection{How much of an architectural difference is real}

The seed analysis answers this directly rather than by inference. Across
four architectures and five seeds, the spread between architectures was 0.66
times the spread from initialisation in the cross-patient setting and 1.09
times it cross-cohort. At that resolution the architectures are not
separable: which one performed best cross-cohort changed with the seed, and
three of the four won at least once across five runs that differed in
nothing else. An architecture ranking computed from a single seed, which is
what the literature reports, is reporting initialisation as much as design.

This is the most direct evidence the study offers on the point. Stated
precisely: among the architectures examined, and at the resolution the
available data permit, architecture contributes less to cross-cohort
variation than either representation choice or initialisation. The claim is
not that architectures are equivalent in general, nor that the remaining ten
would behave the same way, but that the differences observed here in the two
transfer settings are no larger than the variation produced by
initialisation, and so cannot be attributed to the architecture on this
evidence. Only within patients does architecture separate the models by more
than that noise (ratio 2.39), and that is the setting furthest from
deployment. Whether the same holds for the remaining ten architectures, or
at larger sample sizes where smaller differences would become resolvable, is
not settled here.

It also indicates where a benchmark of this kind is most fragile, though
less simply than might be expected. Within-patient evaluation was the most
seed-sensitive setting on average, but the sensitivity was concentrated in
one architecture: EEGNet's mean per-subject seed standard deviation was
0.072 in S1 against ShallowFBCSPNet's 0.009 in the same setting, the latter
lower than its own cross-patient and cross-cohort figures. Reproducibility
is a property of the architecture-and-setting pair rather than of either
alone, and cannot be assumed from the setting. A benchmark reporting a
single seed is not measuring what it claims to.

\subsection{Protocol dependence and comparability}

Published figures for this task span 0.61 AUC to \SI{99.7}{\percent} accuracy.
Our prevalence sweep shows that AUPRC alone moves by a factor of 165 across
the protocols in use, with the underlying discriminative ability unchanged.
Two conclusions follow. First, AUPRC is uninterpretable without the base
rate against which it is to be read, and should be reported alongside it. Second, AUROC, while prevalence-invariant
and therefore stable, cannot reveal the difference and should not be used
alone to argue that two systems are comparable.

The general point is that a metric whose chance level is set by the
composition of the evaluation set cannot be interpreted without that
composition. In seizure detection the composition is a design choice, made
when the evaluation windows are selected, and it is not always reported.

\subsection{Clinical implications}

At the prevalence of continuous monitoring, the precision achievable at
clinically useful sensitivity remains low: AUPRC 0.005 against a base rate of
0.0034, a lift of under 4$\times$. Reported figures obtained on rebalanced
data do not describe this operating point.

The event-level results make the same point in the units a clinical reader
uses, and on continuous recording they can be read directly. A cross-cohort
AUROC of 0.843 corresponds to detecting \SI{22.1}{\percent} of seizures at
one false alarm per hour and \SI{49.8}{\percent} at two. Within patients the
same models reach \SI{77.8}{\percent} at one false alarm per hour with a
median delay of 8 s, which is close to a usable operating point; across
cohorts, at 16 s and a fifth of the sensitivity, it is not.

Two things follow. The gap between window-level and event-level framing is
larger than any difference this study finds between architectures, so a
system reported only at the window level cannot be assessed for clinical
readiness. And the representation gap that AUROC renders as 0.157 appears at
the operating point as a twentyfold difference in sensitivity
(\SI{22.1}{\percent} against \SI{1.1}{\percent} at one false alarm per hour),
which is the scale at which the choice actually matters.

Calibration compounds this. The models were over-confident in every setting,
including the within-patient one, so a threshold chosen during development
will not transfer even to the same patient's later recording without
recalibration.

The per-subject results matter more than the cohort mean for deployment.
Performance on an individual subject ranged from 0.141 to 0.913, and the
failures cannot be anticipated in advance for a new patient. A single global
threshold is therefore not appropriate; per-patient calibration, or an
abstention mechanism that declines to report on subjects where the model is
unreliable, would be more defensible. Establishing which subjects fail, and
why, is a more useful direction than a further increment in mean performance.

\subsection{What the design does and does not establish}

Three features of the design bound what can be concluded from it, and each
bounds a different claim.

\emph{The precision of the estimates.} Resampling subjects rather than
windows gives honest but wide intervals where the held-out set is small: a
mean \SI{95}{\percent} width of 0.208 AUROC across patients, against 0.107
across cohorts. The two generalisation costs, 0.063 and 0.039 in the nested
direction, are therefore estimates with intervals of their own, and the
interval on their ratio, 0.37 to 1.35, includes unity. The comparison is
consequently an ordering supported by its consistency --- 10 of 14
architectures in the nested direction, 12 of 14 in the reverse, and the same
ordering under an independent event-level metric --- rather than a
quantitative claim about how much cheaper the cohort boundary is. A design
with more architectures, or with a third cohort providing further nested
sequences, would be needed to resolve the magnitude.

\emph{The within-patient estimate.} S1 admits eight of the twenty-four
CHB-MIT recording sets, and averaged over architectures it is the most seed-sensitive of
the three settings. The restricted coverage is a property of the data rather
than of the threshold chosen: in several subjects every annotated seizure
falls in the final \SI{30}{\percent} of the recording, so a temporal split
leaves no positive training examples at all. The seed sensitivity is not
uniform across architectures --- it is concentrated in EEGNet, while
ShallowFBCSPNet is more stable in S1 than in either transfer setting --- so
it should be read as a property of particular architecture-and-setting pairs
rather than of S1 as such. Either way the consequence is confined to S1: the
cross-patient and cross-cohort estimates, on which the central comparison
depends, are drawn from the whole cohort.

\emph{The scope of the evidence.} Two cohorts, both focal epilepsy, both
recorded in pre-surgical monitoring units with medication tapered rather
than in the ambulatory conditions an eventual application targets. Siena
contributes 3047 subsampled windows against 43\,198 for CHB-MIT, so the two
transfer directions are not symmetric tests, and Siena also proved the
easier of the two cohorts to fit (leave-one-subject-out 0.747 against 0.718,
SD 0.102 against 0.123). Only the CHB-MIT
$\rightarrow$ Siena direction nests inside the within- and cross-patient
settings, so the reverse direction functions as a replication rather than as
a second estimate of the same quantity; that it points the same way is
supporting evidence, not confirmation. A third cohort would supply further
nested sequences and would test the result properly.

\subsection{Further limitations}

Event-level rates on the full continuous set use a denominator of real
recording time and are comparable both with the handcrafted results and with
the clinical literature. The subsampled event-level figures reported
alongside them, retained for continuity with the sweep, use effective rather
than wall-clock time and are optimistic; they should be read only against
each other.

Hyperparameters were held fixed across architectures, which is what makes
the comparison a comparison, but may disadvantage models that require
tuning. Seed variance was characterised for four of the fourteen
architectures; extending it to the rest would establish whether the figures
reported here are typical of the whole set. The
subject-identity probe is one-vs-rest across 38 subjects while the cohort
probe is binary: both are measured identically, but the tasks are not
equally hard, so that comparison is indicative rather than definitive.

\subsection{Future work}

Three directions follow. Event-level metrics at native prevalence would
establish whether any of these systems reaches a clinically usable operating
point. A third cohort, ideally ambulatory rather than pre-surgical, would test
whether the cohort--patient relationship reported here generalises. And the
per-subject failures, which are large, reproducible across architectures and
currently unexplained, are the most tractable target for improvement: a method
that identified in advance which patients a detector will fail on would be
more clinically valuable than one that raised the mean.

%% =====================================================================
\section{Conclusion}
%% =====================================================================

Under a single protocol applied to two public scalp EEG cohorts, transferring a seizure
detector to a new recording centre was consistently no more costly than
transferring it to a new patient, and generally less so (ratio 0.62,
\SI{95}{\percent} CI 0.37 to 1.35 in the nested direction; 0.32, 0.02 to 0.72
in the reverse). The ordering is what these data establish; the multiple is
not, and in a
substantial fraction of cases it was smaller still. This held across fourteen
architectures spanning three orders of magnitude in parameter count, which
differed among themselves far less than they differed from handcrafted
spectral features. Cohort identity and subject identity were similarly
encoded in the feature space, and the components that shifted most between
cohorts were largely not the discriminative ones.

Under the evaluation protocol used here, cross-cohort performance in these
two cohorts appears to be influenced more strongly by the choice of
representation than by the choice of architecture, and more strongly by
inter-patient variability than by the additional cohort shift measured
between these two cohorts. Separately, a substantial part of the spread in
published performance figures is attributable to evaluation prevalence: AUPRC
moved by two orders of magnitude across the protocols in current use while the
underlying discriminative ability did not change. Reported figures for this
task are not comparable unless the evaluation prevalence is stated and the
evaluation is performed on continuous recording. Scored as a clinician would score
it --- at the event level, on continuous recording, at one false alarm per
hour --- the best of these systems detects around a fifth of seizures in a
cohort it has not seen, and the handcrafted baseline around one in a
hundred.

%% =====================================================================

\section*{Data availability}
Both datasets are public: the CHB-MIT Scalp EEG Database and the Siena Scalp
EEG Database, distributed through PhysioNet. No new human data were collected.

\section*{Code availability}
Analysis code, the window-level manifest and all per-fold results are
available at \url{https://github.com/<user>/eeg-cross-cohort-seizure-detection},
archived at \url{https://doi.org/<Zenodo DOI on release>}.

\section*{Declaration of competing interest}
The author reports no competing interests.

\section*{Funding}
This research received no specific grant from any funding agency.

%% =====================================================================
%%  BIBLIOGRAPHY, EMBEDDED
%%
%%  Held inline rather than generated by BibTeX, so the file compiles with
%%  pdflatex alone and cannot fall out of step with an external .bib.
%%  Entries are in order of first citation. references.bib is kept in the
%%  repository and carries the same entries, for anyone who prefers to
%%  regenerate this list.
%%
%%  To switch back to BibTeX: delete this environment and restore
%%      %%      \bibliography{references}
%% =====================================================================
\begin{thebibliography}{24}
\providecommand{\natexlab}[1]{#1}

\bibitem{shoeb2009}
A.~H. Shoeb.
\newblock \emph{Application of Machine Learning to Epileptic Seizure Onset
  Detection and Treatment}.
\newblock PhD thesis, Massachusetts Institute of Technology, 2009.

\bibitem{hybridconv2025}
P.~K. Upadhyay, P.~Kumar, M.~K. Panda, and A.~K. Samantaray.
\newblock Seamless integration for enhanced seizure prediction using
  HybridConvMobileNet on Typhoon HIL.
\newblock \emph{Scientific Reports}, 15(1):45480, 2025.

\bibitem{jemal2024}
I.~Jemal, L.~Abou-Abbas, K.~Henni, A.~Mitiche, and N.~Mezghani.
\newblock Domain adaptation for {EEG}-based, cross-subject epileptic seizure
  prediction.
\newblock \emph{Frontiers in Neuroinformatics}, 18:1303380, 2024.

\bibitem{crosssubjsurvey2026}
T.~Li, Y.~Yan, F.~Dou, W.~Song, and X.~Zhang.
\newblock Cross-subject generalization for {EEG} decoding: a survey of deep
  learning methods.
\newblock \emph{Progress in Biomedical Engineering}, 8(2):022013, 2026.

\bibitem{saha2024}
E.~Ali, M.~Angelova, and C.~Karmakar.
\newblock Epileptic seizure detection using {CHB-MIT} dataset: the overlooked
  perspectives.
\newblock \emph{Royal Society Open Science}, 11(5):230601, 2024.

\bibitem{saito2015}
T.~Saito and M.~Rehmsmeier.
\newblock The precision-recall plot is more informative than the {ROC} plot
  when evaluating binary classifiers on imbalanced datasets.
\newblock \emph{PLoS ONE}, 10(3):e0118432, 2015.

\bibitem{discoverappl2026}
S.~Ghosh, V.~Singh, and O.~S. Powar.
\newblock Seizure detection using {EEG} on the {CHB-MIT} dataset via
  multi-domain feature engineering and classical machine learning.
\newblock \emph{Discover Applied Sciences}, 8(3):312, 2026.

\bibitem{consumereeg2026}
C.~R. Paredes~Ocaranza, B.~Yun, and E.~D. Paredes~Ocaranza.
\newblock Traditional machine learning outperforms {EEGNet} for consumer-grade
  {EEG} emotion recognition: a comprehensive evaluation with cross-dataset
  validation.
\newblock \emph{Sensors}, 25(23):7262, 2025.

\bibitem{conrad2022}
E.~C. Conrad, J.~M. Bernabei, N.~Sinha, N.~J. Ghosn, J.~M. Stein, R.~T.
  Shinohara, and B.~Litt.
\newblock Addressing spatial bias in intracranial {EEG} functional connectivity
  analyses for epilepsy surgical planning.
\newblock \emph{Journal of Neural Engineering}, 19(5):056019, 2022.

\bibitem{donoghue2020}
T.~Donoghue, M.~Haller, E.~J. Peterson, P.~Varma, P.~Sebastian, R.~Gao,
  T.~Noto, A.~H. Lara, J.~D. Wallis, R.~T. Knight, A.~Shestyuk, and B.~Voytek.
\newblock Parameterizing neural power spectra into periodic and aperiodic
  components.
\newblock \emph{Nature Neuroscience}, 23(12):1655--1665, 2020.

\bibitem{schirrmeister2017}
R.~T. Schirrmeister, J.~T. Springenberg, L.~D.~J. Fiederer, M.~Glasstetter,
  K.~Eggensperger, M.~Tangermann, F.~Hutter, W.~Burgard, and T.~Ball.
\newblock Deep learning with convolutional neural networks for {EEG} decoding
  and visualization.
\newblock \emph{Human Brain Mapping}, 38(11):5391--5420, 2017.

\bibitem{lawhern2018}
V.~J. Lawhern, A.~J. Solon, N.~R. Waytowich, S.~M. Gordon, C.~P. Hung, and
  B.~J. Lance.
\newblock {EEGNet}: a compact convolutional neural network for {EEG}-based
  brain--computer interfaces.
\newblock \emph{Journal of Neural Engineering}, 15(5):056013, 2018.

\bibitem{salami2022}
A.~Salami, J.~Andreu-Perez, and H.~Gillmeister.
\newblock {EEG-ITNet}: an explainable inception temporal convolutional network
  for motor imagery classification.
\newblock \emph{IEEE Access}, 10:36672--36685, 2022.

\bibitem{ingolfsson2020}
T.~M. Ingolfsson, M.~Hersche, X.~Wang, N.~Kobayashi, L.~Cavigelli, and
  L.~Benini.
\newblock {EEG-TCNet}: an accurate temporal convolutional network for embedded
  motor-imagery brain--machine interfaces.
\newblock In \emph{IEEE International Conference on Systems, Man, and
  Cybernetics (SMC)}, pages 2958--2965. IEEE, 2020.

\bibitem{borra2020}
D.~Borra, S.~Fantozzi, and E.~Magosso.
\newblock Interpretable and lightweight convolutional neural network for {EEG}
  decoding: application to movement execution and imagination.
\newblock \emph{Neural Networks}, 129:55--74, 2020.

\bibitem{wei2019}
C.-S. Wei, T.~Koike-Akino, and Y.~Wang.
\newblock Spatial component-wise convolutional network ({SCCNet}) for
  motor-imagery {EEG} classification.
\newblock In \emph{9th International IEEE/EMBS Conference on Neural Engineering
  (NER)}, pages 328--331. IEEE, 2019.

\bibitem{santamaria2020}
E.~Santamar{\'i}a-V{\'a}zquez, V.~Mart{\'i}nez-Cagigal, F.~Vaquerizo-Villar,
  and R.~Hornero.
\newblock {EEG-Inception}: a novel deep convolutional neural network for
  assistive {ERP}-based brain--computer interfaces.
\newblock \emph{IEEE Transactions on Neural Systems and Rehabilitation
  Engineering}, 28(12):2773--2782, 2020.

\bibitem{chen2024}
X.~Chen, X.~Teng, H.~Chen, Y.~Pan, and P.~Geyer.
\newblock Toward reliable signals decoding for electroencephalogram: a
  benchmark study to {EEGNeX}.
\newblock \emph{Biomedical Signal Processing and Control}, 87:105475, 2024.

\bibitem{altaheri2022}
H.~Altaheri, G.~Muhammad, and M.~Alsulaiman.
\newblock Physics-informed attention temporal convolutional network for
  {EEG}-based motor imagery classification.
\newblock \emph{IEEE Transactions on Industrial Informatics}, 19(2):2249--2258,
  2022.

\bibitem{kostas2020}
D.~Kostas and F.~Rudzicz.
\newblock Thinker invariance: enabling deep neural networks for {BCI} across
  more people.
\newblock \emph{Journal of Neural Engineering}, 17(5):056008, 2020.

\bibitem{ouahidi2023}
Y.~El~Ouahidi, V.~Gripon, B.~Pasdeloup, G.~Bouallegue, N.~Farrugia, and
  G.~Lioi.
\newblock A strong and simple deep learning baseline for {BCI} {MI} decoding.
\newblock \emph{arXiv preprint arXiv:2309.07159}, 2023.

\bibitem{song2022}
Y.~Song, Q.~Zheng, B.~Liu, and X.~Gao.
\newblock {EEG} Conformer: convolutional transformer for {EEG} decoding and
  visualization.
\newblock \emph{IEEE Transactions on Neural Systems and Rehabilitation
  Engineering}, 31:710--719, 2022.

\bibitem{jing2023}
J.~Jing, W.~Ge, S.~Hong, M.~B. Fernandes, Z.~Lin, C.~Yang, S.~An, A.~F.
  Struck, A.~Herlopian, I.~Karakis, J.~J. Halford, M.~C. Ng, E.~L. Johnson,
  B.~L. Appavu, R.~A. Sarkis, G.~Osman, P.~W. Kaplan, M.~B. Dhakar, L.~A.
  Jayagopal, Z.~Sheikh, O.~Taraschenko, S.~Schmitt, H.~A. Haider, J.~A. Kim,
  C.~B. Swisher, N.~Gaspard, M.~C. Cervenka, A.~A. Rodriguez~Ruiz, J.~W. Lee,
  M.~Tabaeizadeh, E.~J. Gilmore, K.~Nordstrom, J.~Y. Yoo, M.~G. Holmes, S.~T.
  Herman, J.~A. Williams, J.~Pathmanathan, F.~A. Nascimento, Z.~Fan,
  S.~Nasiri, M.~M. Shafi, S.~S. Cash, D.~B. Hoch, A.~J. Cole, E.~S. Rosenthal,
  S.~F. Zafar, J.~Sun, and M.~B. Westover.
\newblock Development of expert-level classification of seizures and rhythmic
  and periodic patterns during {EEG} interpretation.
\newblock \emph{Neurology}, 100(17):e1750--e1762, 2023.

\bibitem{gramfort2013}
A.~Gramfort, M.~Luessi, E.~Larson, D.~A. Engemann, D.~Strohmeier, C.~Brodbeck,
  R.~Goj, M.~Jas, T.~Brooks, L.~Parkkonen, and M.~H{\"a}m{\"a}l{\"a}inen.
\newblock {MEG} and {EEG} data analysis with {MNE-Python}.
\newblock \emph{Frontiers in Neuroscience}, 7:267, 2013.

\end{thebibliography}

\end{document}

%% =====================================================================
%%  SWITCHING TO IOPART FOR SUBMISSION
%%
%%  1. Download iopart.cls and iopart-num.bst from IOP and upload both.
%%  2. Replace the preamble with:
%%       \documentclass[12pt]{iopart}
%%       \usepackage{iopams}
%%       \usepackage[T1]{fontenc}\usepackage[utf8]{inputenc}
%%       \usepackage{graphicx}\usepackage{amsmath}\usepackage{booktabs}
%%       \usepackage{siunitx}\usepackage[hidelinks]{hyperref}
%%     iopart defines \etal, so delete the \newcommand.
%%     iopart sets its own margins, so delete the geometry line.
%%  3. Replace the title block with \title / \author / \address / \ead /
%%     abstract / \submitto{\JNE} / \maketitle.
%%  4. Replace the bibliography with:
%%       \section*{References}
%%       \bibliographystyle{iopart-num}
%%       \bibliography{references}
%%  5. Change every \citep to \cite.
%% =====================================================================
